\subsection{卷积积的向量表达}

命题 6. --- 设 X、Y、Z 为局部紧空间，\(\varphi\) 为从 \(X \times Y\) 到 Z 的连续映射，\(\lambda, \mu\) 为分别定义在 X、Y 上的测度。\(\lambda\) 与 \(\mu\) 关于 \(\varphi\) 可卷积，当且仅当映射 \((x,y) \mapsto \varepsilon_{\varphi(x,y)} = \varepsilon_x * \varepsilon_y\)，从 \(X \times Y\) 到 \(\mathcal{M}(Z)\)，是标量 \((\lambda \otimes \mu)\)-可积的（相对于拓扑 \(\sigma(\mathcal{M}(Z), \mathcal{K}(Z))\)）；此时

\[
\lambda * \mu = \int_ {\mathrm{X} \times \mathrm{Y}} \left(\varepsilon_ {x} * \varepsilon_ {y}\right) d \lambda (x) d \mu (y).
\]

说 \(\lambda\) 与 \(\mu\) 关于 \(\varphi\) 可卷积，是指对每个 \(f \in \mathcal{H}(\mathbf{Z})\)，\(f \circ \varphi\) 都是 \((\lambda \otimes \mu)\)-可积的；也就是说，对每个 \(f \in \mathcal{H}(\mathbf{Z})\)，函数 \((x,y) \mapsto \langle f, \varepsilon_{\varphi(x,y)} \rangle\) 是 \((\lambda \otimes \mu)\)-可积的；换言之，映射 \((x,y) \mapsto \varepsilon_{\varphi(x,y)}\) 从 \(\mathrm{X} \times \mathrm{Y}\) 到 \(\mathcal{M}(\mathbf{Z})\) 是标量 \((\lambda \otimes \mu)\)-可积的（相对于 \(\sigma (\mathcal{M}(\mathbb{Z}),\mathcal{K}(\mathbb{Z}))\)）。若此成立，则

\[
\langle \lambda * \mu , f \rangle = \int f (\varphi (x, y)) d \lambda (x) d \mu (y) = \int_ {\mathrm{X} \times \mathrm{Y}} \left\langle \varepsilon_ {\varphi (x, y)}, f \right\rangle d \lambda (x) d \mu (y),
\]

从而 \(\lambda * \mu = \int_{\mathbf{X} \times \mathbf{Y}} \varepsilon_{\varphi(x, y)} d\lambda(x) d\mu(y)\)。

命题 7. --- 设 X、Y、Z 为局部紧空间，\(\varphi\) 为从 \(X \times Y\) 到 Z 的连续映射，\(\lambda, \mu\) 为分别定义在 X、Y 上的测度。假设对每个 \(x \in X\)，\(\varepsilon_x\) 与 \(\mu\) 关于 \(\varphi\) 可卷积。\(\lambda\) 与 \(\mu\) 关于 \(\varphi\) 可卷积，当且仅当映射 \(x \mapsto \varepsilon_x * |\mu|\) 从 X 到 \(\mathcal{M}(Z)\) 是标量 \(\lambda\)-可积的（相对于拓扑 \(\sigma(\mathcal{M}(Z), \mathcal{K}(Z))\)）；此时 \(\lambda * \mu = \int_{X} (\varepsilon_x * \mu) d\lambda(x)\)。

假设 \(\lambda\) 与 \(\mu\) 关于 \(\varphi\) 可卷积。对每个 \(f \in \mathcal{K}(\mathbb{Z})\)，\(f \circ \varphi\) 是 \((|\lambda| \otimes |\mu|)\)-可积的，因此函数 \(x \mapsto \int_{\mathbb{Y}} f(\varphi(x, y)) d|\mu|(y) = \langle f, \varepsilon_x * |\mu| \rangle\)（根据假设，它对所有 \(x \in \mathbb{X}\) 有定义）是 \(\lambda\)-可积的；于是 \(x \mapsto \varepsilon_x * |\mu|\) 是标量 \(\lambda\)-可积的（相对于 \(\sigma(\mathcal{M}(\mathbb{Z}), \mathcal{K}(\mathbb{Z}))\)），并且

\[
\langle f, \lambda * \mu \rangle = \int_ {\mathrm{X}} d \lambda (x) \int_ {\mathrm{Y}} f (\varphi (x, y)) d \mu (y) = \int_ {\mathrm{X}} \left\langle f, \varepsilon_ {x} * \mu \right\rangle d \lambda (x),
\]

从而 \(\lambda * \mu = \int_{\mathbf{X}} (\varepsilon_x * \mu) d\lambda(x)\)。反之，假设映射 \(x \mapsto \varepsilon_x * |\mu|\) 从 X 到 \(\mathcal{M}(\mathbf{Z})\) 是标量 \(\lambda\)-可积的（相对于 \(\sigma(\mathcal{M}(\mathbf{Z}), \mathcal{K}(\mathbf{Z}))\)）。令 \(f \in \mathcal{H}_+(\mathbf{Z})\)。则函数 \((x,y) \mapsto f(\varphi(x,y))\) 连续，并且（Ch. V, §8, No.~3, 命题 5）

\[
\begin{array}{r l} \iint^ {*} f (\varphi (x, y)) d | \lambda | (x) d | \mu | (y) & = \int^ {*} d | \lambda | (x) \int^ {*} f (\varphi (x, y)) d | \mu | (y) \\ & = \int^ {*} \langle f, \varepsilon_ {x} * | \mu | \rangle d | \lambda | (x) <   + \infty . \end{array}
\]

因此 \(f \circ \varphi\) 是 \((\lambda \otimes \mu)\)-可积的，从而 \(\lambda\) 与 \(\mu\) 关于 \(\varphi\) 可卷积。

